\documentclass[10pt,a4paper]{amsart}
\usepackage[margin=19mm]{geometry}
\usepackage{amsmath,amssymb,mathtools}
\usepackage[scr=rsfs,cal=euler]{mathalfa}
\usepackage{xfrac}
\usepackage[unicode,hidelinks,bookmarks=false]{hyperref}
\newcommand{\ie}{i.e.\ }
\newcommand{\cf}{cf.\ }
\newcommand{\resp}{respectively\ }
\newcommand{\Subsection}{\subsection}
\providecommand{\nobreakdash}{\nobreak\hbox{-}}
\setlength{\emergencystretch}{2em}
\multlinegap0pt
\usepackage{bm}
\usepackage{bbm}
\usepackage{dsfont}
\usepackage{xspace}
\usepackage{cancel}
\usepackage{mathtools}
\usepackage{amsthm}
\makeatletter
\@ifundefined{theorem}{\newtheorem{theorem}{Theorem}}{}
\@ifundefined{lemma}{\newtheorem{lemma}[theorem]{Lemma}}{}
\makeatother
\title[Vector-Valued Period Polynomials and Zeta Values of Quadrati]{Vector-Valued Period Polynomials and Zeta Values of Quadratic Fields}
\author{Yeong-Wook Kwon, Subong Lim,, Wissam Raji}
\date{Source: arXiv:2602.00572v1 (CC BY 4.0)}
\begin{document}
\maketitle

\noindent\emph{Source and licence.} Vector-Valued Period Polynomials and Zeta Values of Quadratic Fields. Version arXiv:2602.00572v1 (CC BY 4.0), distributed under the Creative Commons Attribution 4.0 International licence, as recorded in the arXiv licence metadata for this exact version and shown on the abstract page https://arxiv.org/abs/2602.00572v1. Full rights notice: licenses/arxiv-2602.00572-CC-BY-4.0.txt.\par\smallskip

\noindent\textit{Editorial scope.} Counted unit: the Lemma whose source label is \texttt{lem-pairtoideal} in the article (source0.tex lines 811--825, source label \texttt{lem-pairtoideal}), delivered together with the article's own complete original proof of it (source0.tex lines 827--871, one \texttt{proof} environment of 45 lines). No mathematics is written, repaired or re-derived: statement and proof are the article's own text line for line. Required context retained uncounted: none. Every definition, display and label of those lines is retained unchanged. Omitted from this package and not redistributed: 1--810, 826--826, 872--1245. Nothing inside a retained excerpt is omitted or altered: every retained body is delivered exactly as the article's lines are written, so no omission receipt of any kind accompanies this package. Citations: the retained excerpts carry no citation command at all, so no bibliography entry of the article is transcribed and no reference is redistributed. Reference adaptation: none was needed and none was made; every reference inside the delivered unit resolves inside it, so no unresolved reference or citation is produced. Printed slips kept literal: no printed word, symbol, hypothesis or number of the counted statement or of its proof was altered. Layout and class: the article's own class is \texttt{amsart} with packages including amsmath, amsfonts, amsthm, amssymb, xy, enumitem, tikz-cd, hyperref, mathrsfs, tikz, babel, graphicx, centernot, mathtools; the wrapper is the lane's pinned \texttt{amsart} editorial wrapper at 10pt on A4, which restates the theorem-like environments the retained text needs and adds the article's own macro definitions that the retained text needs (none), with extra wrapper packages (bm, bbm, dsfont, xspace, cancel, mathtools). The delivered numbering is the unit's own. Assets: no retained span contains a figure environment, a graphics-inclusion command, a PDF-page inclusion, a table or a TikZ picture; no image file is copied into this package, the package declares no asset dependency and no image file is redistributed with it. Licence: version arXiv:2602.00572v1 of this article is distributed under the Creative Commons Attribution 4.0 International licence, as recorded in the arXiv licence metadata for this exact version and shown on the abstract page https://arxiv.org/abs/2602.00572v1. A full rights notice is delivered at licenses/arxiv-2602.00572-CC-BY-4.0.txt. Characters: every retained excerpt is pure ASCII, so the delivered engine, XeLaTeX with its default Latin Modern text font, reports no missing character.
\begin{lemma}\label{lem-pairtoideal}
Let $Q=[A,B,C]\in\mathcal{Q}_{D}^{0}$.
For $(u,v)\in\mathbb{Z}^{2}$, define
$$\mathfrak{b}_{u,v}(Q):=(u-v\alpha_{Q}')\mathfrak{a}_{Q}.$$
Then
\begin{enumerate}
    \item[\textnormal{(i)}] If $Q(u,v)>0$, then $\mathfrak{b}_{u,v}(Q)$ is an integral ideal of $\mathcal{O}_{K}$.
    \item[\textnormal{(ii)}] One has $\mathrm{Nm}(\mathfrak{b}_{u,v}(Q))=Q(u,v)$.
    \item[\textnormal{(iii)}] The map $(u,v)\mapsto\mathfrak{b}_{u,v}(Q)$ induces a well-defined bijection between
    \begin{equation}\label{eq-pairideal}
        \{(u,v)\in\mathbb{Z}^2/\mathrm{SL}_{2}(\mathbb{Z})_{Q}:\ Q(u,v)>0\}\quad\text{and}\quad
        \{\mathfrak{b}\subset\mathcal{O}_{K}:[\mathfrak{b}]=[\mathfrak{a}_{Q}]~\text{in}~\mathrm{Cl}^{+}(\mathcal{O}_{K})\}.
    \end{equation}
\end{enumerate}
\end{lemma}

\begin{proof}
Note that
$$A\alpha_{Q}'=\omega_{B}'=\frac{-B-\sqrt{D}}{2}\quad\text{and}\quad\alpha_{Q}'\omega_{B}=\frac{(-B-\sqrt{D})(-B+\sqrt{D})}{4A}=\frac{B^{2}-D}{4A}=C\in\mathbb{Z}.$$

(i) We compute
$$\mathfrak{b}_{u,v}(Q)=(u-v\alpha_{Q}')(A,\omega_{B})=(uA-vA\alpha_{Q}',u\omega_{B}-v\alpha_{Q}'\omega_{B})=(uA-v\omega_{B}',u\omega_{B}-vC).$$
Since $\omega_{B},\omega_{B}'\in\mathcal{O}_{K}$, we conclude that $\mathfrak{b}_{u,v}(Q)$ is an integral ideal.

(ii) Since $\alpha_Q$ and $\alpha_Q'$ are conjugates, we have
$$\mathrm{Nm}(u-v\alpha_{Q}')=(u-v\alpha_{Q}')(u-v\alpha_{Q}).$$
Using $\alpha_{Q}+\alpha_{Q}'=-B/A$ and $\alpha_{Q}\alpha_{Q}'=C/A$, we obtain
$$\mathrm{Nm}(u-v\alpha_{Q}')=u^{2}-uv(\alpha_{Q}+\alpha_{Q}')+v^{2}\alpha_{Q}\alpha_{Q}'=u^{2}+\frac{B}{A}uv+\frac{C}{A}v^{2}=\frac{Q(u,v)}{A}.$$
Moreover, $\mathrm{Nm}(\mathfrak{a}_{Q})=A$. Therefore,
$$\mathrm{Nm}(\mathfrak{b}_{u,v}(Q))=\mathrm{Nm}(u-v\alpha_{Q}')\cdot\mathrm{Nm}(\mathfrak{a}_{Q})=Q(u,v).$$

(iii) If $Q(u,v)>0$, then $\mathrm{Nm}(u-v\alpha_Q')=Q(u,v)/A>0$.
In a real quadratic field, a principal ideal generated by an element of positive norm is narrow principal (replace the generator by its negative if necessary).
Hence, $[(u-v\alpha_{Q}')\mathfrak{a}_{Q}]=[\mathfrak{a}_{Q}]$ in $\mathrm{Cl}^{+}(\mathcal{O}_{K})$.

We next check that the map is well-defined on $\mathbb{Z}^{2}/\mathrm{SL}_{2}(\mathbb{Z})_{Q}$.
Let $M=\left(\begin{smallmatrix} \alpha & \beta \\ \gamma & \delta\end{smallmatrix}\right)\in\mathrm{SL}_{2}(\mathbb{Z})_{Q}$ and set $(u',v'):=(u,v)M$.
Then $M$ fixes both $\alpha_{Q}$ and $\alpha_{Q}'$ under the fractional linear transformation, and $\mathrm{Nm}(\epsilon_{M})=1$, where $\epsilon_{M}=\gamma\alpha_{Q}'+\delta$.
Moreover, one has
$$u'-v'\alpha_{Q}'=\frac{u-v\alpha_{Q}'}{\gamma\alpha_{Q}'+\delta}=\frac{u-v\alpha_{Q}'}{\epsilon_{M}}.$$
Since $(\epsilon_{M})=\mathcal{O}_{K}$,
$$\mathfrak{b}_{u',v'}(Q)=(u'-v'\alpha_{Q}')\mathfrak{a}_{Q}=(u-v\alpha_{Q}')\mathfrak{a}_{Q}=\mathfrak{b}_{u,v}(Q).$$
Thus, the map is well-defined on the quotient.

For surjectivity, let $\mathfrak{b}\subset\mathcal{O}_{K}$ satisfy $[\mathfrak{b}]=[\mathfrak{a}_{Q}]$ in $\mathrm{Cl}^{+}(\mathcal{O}_{K})$.
Then there exists totally positive $\beta\in K^{\times}$ with $\mathfrak{b}=(\beta)\mathfrak{a}_{Q}$.
Since $\mathfrak{b}\subset\mathcal{O}_{K}$, we have $\beta\in\mathfrak{a}_{Q}^{-1}$.
One checks that
$$\mathfrak{a}_{Q}^{-}=\mathbb{Z}+\mathbb{Z}\alpha_{Q}'.$$
Indeed, $1\in\mathfrak{a}_{Q}^{-1}$ and $\alpha_{Q}'\in\mathfrak{a}_{Q}^{-1}$ since $\alpha_Q'A=\omega_{B}'\in\mathcal{O}_{K}$ and $\alpha_Q'\omega_B=C\in\mathbb{Z}$.
$$(A,\omega_{B})(\mathbb{Z}+\mathbb{Z}\alpha_{Q}')=(A,\omega_{B},A\alpha_{Q}',\omega_{B}\alpha_{Q}')=(A,\omega_{B},\omega_{B}',C)$$
contains $B=-(\omega_{B}+\omega_{B}')\in\mathbb{Z}$, hence contains $(A,B,C)=\mathcal{O}_{K}$ because $\gcd(A,B,C)=1$.
Therefore, $\mathfrak{a}_{Q}^{-1}=\mathbb{Z}+\mathbb{Z}\alpha_{Q}'$. This shows that $\beta=u-v\alpha_{Q}'$ for some $u,v\in\mathbb{Z}$, and $\mathfrak{b}=\mathfrak{b}_{u,v}(Q)$.
Finally, by (ii), 
$$0<\mathrm{Nm}(\mathfrak b)=\mathrm{Nm}(\mathfrak{b}_{u,v}(Q))=Q(u,v).$$

Injectivity follows similarly: if $\mathfrak b_{u,v}(Q)=\mathfrak b_{u',v'}(Q)$, then
$$\frac{u-v\alpha_{Q}'}{u'-v'\alpha_{Q}'}\in\mathcal{O}_{K}'\quad\text{and}\quad\mathrm{Nm}\left(\frac{u-v\alpha_{Q}'}{u'-v'\alpha_{Q}'}\right)=.$$
Such a unit corresponds to an element of $\mathrm{PSL}_{2}(\mathbb{Z})_{Q}$; if necessary, composing with $-I\in\mathrm{SL}_{2}(\mathbb{Z})_{Q}$ gives an element of $\mathrm{SL}_2(\mathbb{Z})_{Q}$.
Hence, $(u,v)$ and $(u',v')$ represent the same class in $\mathbb{Z}^{2}/\mathrm{SL}_{2}(\mathbb{Z})_{Q}$.
\end{proof}

\end{document}
